\section{Projections coordonnées de l'état HMT : équivalence algébrique des écrans et réalisation physique en dimension 11}
\label{sec:cierre-equivalencias-moonshine-teoria-m-20260825}

Soit \(\mathcal X_{\mathrm{HMT}}^{\mathrm{enr}}\) le domaine qui conserve
la phase, l'orientation, le quotient, la retenue, la frontière, l'incidence, la graduation et la
mémoire. Ses composantes exceptionnelle et d'écran sont extraites par les
projections canoniques
\[
 \varepsilon_{\mathrm{exc}}:
 \mathcal X_{\mathrm{HMT}}^{\mathrm{enr}}
 \longrightarrow\mathscr X_{\mathrm{exc}},
 \qquad
 \varepsilon_{\mathrm{scr}}:
 \mathcal X_{\mathrm{HMT}}^{\mathrm{enr}}
 \longrightarrow\mathscr S_{\mathrm{HMT}}.
\]
La première alimente la construction déjà démontrée
\(\mathcal P_{\mathrm{exc}}:\mathscr X_{\mathrm{exc}}\to V^\natural\).
La seconde alimentera le morphisme algébrique
\(\mathcal R_M^{\mathrm{alg}}\) construit plus bas. La réalisation conjointe
est
\[
 \mathcal P_{\mathrm{joint}}^M(x)
 =\left(
   \mathcal P_{\mathrm{exc}}\varepsilon_{\mathrm{exc}}(x),
   \mathcal R_M^{\mathrm{alg}}\varepsilon_{\mathrm{scr}}(x)
  \right).
\]
Cette application conserve l'antécédent commun et maintient les codomaines distincts.

\subsection{Équivalence exacte des écrans réticulaires et conditions de la réalisation undécadimensionnelle}

L'écran réel part de
\[
 V_{12}\cong V_1\oplus V_3\oplus V_{3'}\oplus V_5,
 \qquad
 P_{11}=I_{12}-\frac1{12}J_{12}.
\]
L'état \(K\) produit \(u_K=P_3P_{11}K\ne0\), la droite
\(\ell_K=\mathbb Ru_K\), et
\[
 V_{11}=\ell_K\oplus V_{10},
 \qquad \dim V_{11}=11,
 \qquad \dim V_{10}=10.
\]
Le passage \(12\to11\) élimine le mode uniforme et le passage \(11\to10\) élimine
la droite polarisée par \(K\).

L'écran réticulaire appartient à une autre catégorie. Pour
\[
 L_{26}=\Lambda_{24}\oplus U\cong II_{25,1}
\]
et un vecteur isotrope primitif \(e_+\in U\),
\[
 e_+^\perp=\Lambda_{24}\oplus\mathbb Ze_+,
 \qquad
 \boxed{e_+^\perp/\mathbb Ze_+\cong\Lambda_{24}.}
\]
L'orthogonalité et le quotient éliminent deux directions. Ce théorème n'est pas
une répétition de \(11\to10\) ni une identification automatique à un
espace--temps physique.

\subsection{Morphisme algébrique d'écran et équivalence quotientée}

Soit
\[
 P_{10}:=P_{11}-\frac{u_Ku_K^{\mathsf T}}{\lVert u_K\rVert^2},
 \qquad
 \mathscr G_K:=
 \{(v_{11},v_{10})\in V_{11}\times V_{10}:
             v_{10}=P_{10}v_{11}\}.
\]
La décomposition
\(e_+^\perp=\Lambda_{24}\oplus\mathbb Ze_+\) définit
\[
 q_{24}:e_+^\perp\longrightarrow\Lambda_{24},
 \qquad q_{24}(\lambda+ae_+)=\lambda.
\]
Prenons
\[
 \mathscr S_{\mathrm{HMT}}
 :=V_{12}\times\mathscr D_0\times e_+^\perp,
 \qquad
 \mathscr S_M^{\mathrm{alg}}
 :=\mathscr G_K\times\mathscr D_0\times\Lambda_{24}.
\]

\begin{definition}[Morphisme algébrique d'écran M]
\label{def:cierre-vtm-rm-algebraico}
Pour \(d=(m,w,R_0)\in\mathscr D_0\), nous définissons
\[
 \boxed{
 \mathcal R_M^{\mathrm{alg}}(v,d,y)
 :=\bigl((P_{11}v,P_{10}v),d,q_{24}(y)\bigr).}
\]
\end{definition}

\begin{theorem}[Équivalence algébrique des écrans quotientés]
\label{thm:cierre-vtm-rm-algebraico}
Le morphisme précédent induit un isomorphisme
\[
 \boxed{
 \overline{\mathcal R}_M^{\mathrm{alg}}:
 (V_{12}/V_1)\times\mathscr D_0
 \times(e_+^\perp/\mathbb Ze_+)
 \xrightarrow{\ \cong\ }
 \mathscr S_M^{\mathrm{alg}}.}
\]
Si
\[
 \tau_{\mathrm{src}}([v],d,[y])
 =([v],\mathsf T_0d,[y])
\]
et
\[
 \tau_{\mathrm{tar}}((v_{11},v_{10}),d,\lambda)
 =((v_{11},v_{10}),\mathsf T_0d,\lambda),
\]
alors
\[
 \overline{\mathcal R}_M^{\mathrm{alg}}\tau_{\mathrm{src}}
 =\tau_{\mathrm{tar}}\overline{\mathcal R}_M^{\mathrm{alg}}.
\]
L'énergie compacte \(E_{R_0}(m,w)\) est conservée sous cet
entrelacement et sous des rayons réciproques.
\end{theorem}

\begin{proof}
Le projecteur \(P_{11}\) a pour noyau \(V_1\) et pour image \(V_{11}\). Comme
\(P_{10}P_{11}=P_{10}\), l'application
\[
 [v]\longmapsto(P_{11}v,P_{10}v)
\]
est un isomorphisme de \(V_{12}/V_1\) sur le graphe
\(\mathscr G_K\). L'application \(q_{24}\) est surjective et son noyau est
\(\mathbb Ze_+\), de sorte qu'elle induit l'isomorphisme
\(e_+^\perp/\mathbb Ze_+\cong\Lambda_{24}\). Le facteur
\(\mathscr D_0\) se transporte par l'identité. Le produit de ces trois
isomorphismes est \(\overline{\mathcal R}_M^{\mathrm{alg}}\).

Les projecteurs et \(q_{24}\) n'agissent pas sur \(d\), tandis que les deux
involutions n'agissent que par \(\mathsf T_0\). Le carré commute
par substitution directe. L'invariance énergétique est le
\cref{thm:cierre-vtm-tres-formulaciones-t}.
\end{proof}

Ainsi est construite l'application mathématique de réduction, et pas seulement une
correspondance de rangs. L'écran \(11\to10\), le quotient
\(26\to24\) et la dualité de rayon apparaissent comme des facteurs distincts du
même morphisme.

\subsection{Reconnaissance physique undécadimensionnelle}

L'interprétation comme théorie M est séparée par une application de reconnaissance
\[
 \Xi_M:\mathscr S_M^{\mathrm{alg}}
 \longrightarrow\mathcal X_{11}^{\mathrm{phys}},
\]
dont le codomaine porte au moins
\[
 g_{MN},\qquad C_{MNP},\qquad\psi_M,
\]
l'action et les transformations de supersymétrie. Une complétion
opératoire comprend
\[
 \mathfrak S_M=(\mathcal H_{\mathrm{phys}},H,Q),
 \qquad Q^2=H,
\]
et une isométrie
\[
 W:\mathcal H_{\mathrm{HMT}}\longrightarrow\mathcal H_{\mathrm{phys}}
\]
qui, sur leurs domaines communs, satisfait
\[
 W\widehat H=HW,
 \qquad W\widehat Q=QW,
 \qquad \widehat Q^{\,2}=\widehat H.
\]

\begin{theorem}[Critère de reconnaissance physique compatible]
\label{thm:cierre-vtm-reconocimiento-fisico-m}
Si \(\Xi_M\) préserve le produit scalaire, transporte la direction
\(\ell_K\) vers la direction compacte, entrelace la dualité et satisfait les
identités opératoires précédentes, alors le spectre physique compact est
invariant sous des rayons réciproques et la supercharge transportée satisfait
\(Q^2=H\) sur l'image de \(W\).
\end{theorem}

\begin{proof}
L'entrelacement de \(\Xi_M\) avec \(\tau_{\mathrm{tar}}\), joint au
\cref{thm:cierre-vtm-rm-algebraico}, transporte l'involution unitaire au
secteur compact. Pour la supercharge,
\[
 Q^2W=W\widehat Q^{\,2}=W\widehat H=HW.
\]
La préservation du produit scalaire transforme l'opérateur transporté en
une équivalence unitaire sur l'image.
\end{proof}

L'isomorphisme
\(\overline{\mathcal R}_M^{\mathrm{alg}}\) est mathématique et démontré.
La reconnaissance par \(\Xi_M\) est une interface physique conditionnée par les
champs, l'action, les contraintes de jauge, la supersymétrie, les
oscillateurs, les amplitudes, le contrôle des anomalies et la limite de basse
énergie de la réalisation choisie. Le codomaine physique n'est pas identifié à
\(V^\natural\) : tous deux reçoivent une information coordonnée du même antécédent
par des foncteurs distincts.

Aucune action littérale de \(\mathbb M\) sur les chiffres n'est postulée. Cette phrase est
le contrôle de typage placé à la fin ; elle ne remplace pas la chaîne positive qui
atteint l'action profonde de \(\mathbb M\) dans \(V^\natural\).
